% Options for packages loaded elsewhere
% Options for packages loaded elsewhere
\PassOptionsToPackage{unicode}{hyperref}
\PassOptionsToPackage{hyphens}{url}
\PassOptionsToPackage{dvipsnames,svgnames,x11names}{xcolor}
%
\documentclass[
  letterpaper,
  DIV=11,
  numbers=noendperiod]{scrreprt}
\usepackage{xcolor}
\usepackage{amsmath,amssymb}
\setcounter{secnumdepth}{5}
\usepackage{iftex}
\ifPDFTeX
  \usepackage[T1]{fontenc}
  \usepackage[utf8]{inputenc}
  \usepackage{textcomp} % provide euro and other symbols
\else % if luatex or xetex
  \usepackage{unicode-math} % this also loads fontspec
  \defaultfontfeatures{Scale=MatchLowercase}
  \defaultfontfeatures[\rmfamily]{Ligatures=TeX,Scale=1}
\fi
\usepackage{lmodern}
\ifPDFTeX\else
  % xetex/luatex font selection
\fi
% Use upquote if available, for straight quotes in verbatim environments
\IfFileExists{upquote.sty}{\usepackage{upquote}}{}
\IfFileExists{microtype.sty}{% use microtype if available
  \usepackage[]{microtype}
  \UseMicrotypeSet[protrusion]{basicmath} % disable protrusion for tt fonts
}{}
\makeatletter
\@ifundefined{KOMAClassName}{% if non-KOMA class
  \IfFileExists{parskip.sty}{%
    \usepackage{parskip}
  }{% else
    \setlength{\parindent}{0pt}
    \setlength{\parskip}{6pt plus 2pt minus 1pt}}
}{% if KOMA class
  \KOMAoptions{parskip=half}}
\makeatother
% Make \paragraph and \subparagraph free-standing
\makeatletter
\ifx\paragraph\undefined\else
  \let\oldparagraph\paragraph
  \renewcommand{\paragraph}{
    \@ifstar
      \xxxParagraphStar
      \xxxParagraphNoStar
  }
  \newcommand{\xxxParagraphStar}[1]{\oldparagraph*{#1}\mbox{}}
  \newcommand{\xxxParagraphNoStar}[1]{\oldparagraph{#1}\mbox{}}
\fi
\ifx\subparagraph\undefined\else
  \let\oldsubparagraph\subparagraph
  \renewcommand{\subparagraph}{
    \@ifstar
      \xxxSubParagraphStar
      \xxxSubParagraphNoStar
  }
  \newcommand{\xxxSubParagraphStar}[1]{\oldsubparagraph*{#1}\mbox{}}
  \newcommand{\xxxSubParagraphNoStar}[1]{\oldsubparagraph{#1}\mbox{}}
\fi
\makeatother

\usepackage{color}
\usepackage{fancyvrb}
\newcommand{\VerbBar}{|}
\newcommand{\VERB}{\Verb[commandchars=\\\{\}]}
\DefineVerbatimEnvironment{Highlighting}{Verbatim}{commandchars=\\\{\}}
% Add ',fontsize=\small' for more characters per line
\usepackage{framed}
\definecolor{shadecolor}{RGB}{241,243,245}
\newenvironment{Shaded}{\begin{snugshade}}{\end{snugshade}}
\newcommand{\AlertTok}[1]{\textcolor[rgb]{0.68,0.00,0.00}{#1}}
\newcommand{\AnnotationTok}[1]{\textcolor[rgb]{0.37,0.37,0.37}{#1}}
\newcommand{\AttributeTok}[1]{\textcolor[rgb]{0.40,0.45,0.13}{#1}}
\newcommand{\BaseNTok}[1]{\textcolor[rgb]{0.68,0.00,0.00}{#1}}
\newcommand{\BuiltInTok}[1]{\textcolor[rgb]{0.00,0.23,0.31}{#1}}
\newcommand{\CharTok}[1]{\textcolor[rgb]{0.13,0.47,0.30}{#1}}
\newcommand{\CommentTok}[1]{\textcolor[rgb]{0.37,0.37,0.37}{#1}}
\newcommand{\CommentVarTok}[1]{\textcolor[rgb]{0.37,0.37,0.37}{\textit{#1}}}
\newcommand{\ConstantTok}[1]{\textcolor[rgb]{0.56,0.35,0.01}{#1}}
\newcommand{\ControlFlowTok}[1]{\textcolor[rgb]{0.00,0.23,0.31}{\textbf{#1}}}
\newcommand{\DataTypeTok}[1]{\textcolor[rgb]{0.68,0.00,0.00}{#1}}
\newcommand{\DecValTok}[1]{\textcolor[rgb]{0.68,0.00,0.00}{#1}}
\newcommand{\DocumentationTok}[1]{\textcolor[rgb]{0.37,0.37,0.37}{\textit{#1}}}
\newcommand{\ErrorTok}[1]{\textcolor[rgb]{0.68,0.00,0.00}{#1}}
\newcommand{\ExtensionTok}[1]{\textcolor[rgb]{0.00,0.23,0.31}{#1}}
\newcommand{\FloatTok}[1]{\textcolor[rgb]{0.68,0.00,0.00}{#1}}
\newcommand{\FunctionTok}[1]{\textcolor[rgb]{0.28,0.35,0.67}{#1}}
\newcommand{\ImportTok}[1]{\textcolor[rgb]{0.00,0.46,0.62}{#1}}
\newcommand{\InformationTok}[1]{\textcolor[rgb]{0.37,0.37,0.37}{#1}}
\newcommand{\KeywordTok}[1]{\textcolor[rgb]{0.00,0.23,0.31}{\textbf{#1}}}
\newcommand{\NormalTok}[1]{\textcolor[rgb]{0.00,0.23,0.31}{#1}}
\newcommand{\OperatorTok}[1]{\textcolor[rgb]{0.37,0.37,0.37}{#1}}
\newcommand{\OtherTok}[1]{\textcolor[rgb]{0.00,0.23,0.31}{#1}}
\newcommand{\PreprocessorTok}[1]{\textcolor[rgb]{0.68,0.00,0.00}{#1}}
\newcommand{\RegionMarkerTok}[1]{\textcolor[rgb]{0.00,0.23,0.31}{#1}}
\newcommand{\SpecialCharTok}[1]{\textcolor[rgb]{0.37,0.37,0.37}{#1}}
\newcommand{\SpecialStringTok}[1]{\textcolor[rgb]{0.13,0.47,0.30}{#1}}
\newcommand{\StringTok}[1]{\textcolor[rgb]{0.13,0.47,0.30}{#1}}
\newcommand{\VariableTok}[1]{\textcolor[rgb]{0.07,0.07,0.07}{#1}}
\newcommand{\VerbatimStringTok}[1]{\textcolor[rgb]{0.13,0.47,0.30}{#1}}
\newcommand{\WarningTok}[1]{\textcolor[rgb]{0.37,0.37,0.37}{\textit{#1}}}

\usepackage{longtable,booktabs,array}
\newcounter{none} % for unnumbered tables
\usepackage{calc} % for calculating minipage widths
% Correct order of tables after \paragraph or \subparagraph
\usepackage{etoolbox}
\makeatletter
\patchcmd\longtable{\par}{\if@noskipsec\mbox{}\fi\par}{}{}
\makeatother
% Allow footnotes in longtable head/foot
\IfFileExists{footnotehyper.sty}{\usepackage{footnotehyper}}{\usepackage{footnote}}
\makesavenoteenv{longtable}
\usepackage{graphicx}
\makeatletter
\newsavebox\pandoc@box
\newcommand*\pandocbounded[1]{% scales image to fit in text height/width
  \sbox\pandoc@box{#1}%
  \Gscale@div\@tempa{\textheight}{\dimexpr\ht\pandoc@box+\dp\pandoc@box\relax}%
  \Gscale@div\@tempb{\linewidth}{\wd\pandoc@box}%
  \ifdim\@tempb\p@<\@tempa\p@\let\@tempa\@tempb\fi% select the smaller of both
  \ifdim\@tempa\p@<\p@\scalebox{\@tempa}{\usebox\pandoc@box}%
  \else\usebox{\pandoc@box}%
  \fi%
}
% Set default figure placement to htbp
\def\fps@figure{htbp}
\makeatother


% definitions for citeproc citations
\NewDocumentCommand\citeproctext{}{}
\NewDocumentCommand\citeproc{mm}{%
  \begingroup\def\citeproctext{#2}\cite{#1}\endgroup}
\makeatletter
 % allow citations to break across lines
 \let\@cite@ofmt\@firstofone
 % avoid brackets around text for \cite:
 \def\@biblabel#1{}
 \def\@cite#1#2{{#1\if@tempswa , #2\fi}}
\makeatother
\newlength{\cslhangindent}
\setlength{\cslhangindent}{1.5em}
\newlength{\csllabelwidth}
\setlength{\csllabelwidth}{3em}
\newenvironment{CSLReferences}[2] % #1 hanging-indent, #2 entry-spacing
 {\begin{list}{}{%
  \setlength{\itemindent}{0pt}
  \setlength{\leftmargin}{0pt}
  \setlength{\parsep}{0pt}
  % turn on hanging indent if param 1 is 1
  \ifodd #1
   \setlength{\leftmargin}{\cslhangindent}
   \setlength{\itemindent}{-1\cslhangindent}
  \fi
  % set entry spacing
  \setlength{\itemsep}{#2\baselineskip}}}
 {\end{list}}
\usepackage{calc}
\newcommand{\CSLBlock}[1]{\hfill\break\parbox[t]{\linewidth}{\strut\ignorespaces#1\strut}}
\newcommand{\CSLLeftMargin}[1]{\parbox[t]{\csllabelwidth}{\strut#1\strut}}
\newcommand{\CSLRightInline}[1]{\parbox[t]{\linewidth - \csllabelwidth}{\strut#1\strut}}
\newcommand{\CSLIndent}[1]{\hspace{\cslhangindent}#1}



\setlength{\emergencystretch}{3em} % prevent overfull lines

\providecommand{\tightlist}{%
  \setlength{\itemsep}{0pt}\setlength{\parskip}{0pt}}



 


\KOMAoption{captions}{tableheading}
\makeatletter
\@ifpackageloaded{bookmark}{}{\usepackage{bookmark}}
\makeatother
\makeatletter
\@ifpackageloaded{caption}{}{\usepackage{caption}}
\AtBeginDocument{%
\ifdefined\contentsname
  \renewcommand*\contentsname{Table of contents}
\else
  \newcommand\contentsname{Table of contents}
\fi
\ifdefined\listfigurename
  \renewcommand*\listfigurename{List of Figures}
\else
  \newcommand\listfigurename{List of Figures}
\fi
\ifdefined\listtablename
  \renewcommand*\listtablename{List of Tables}
\else
  \newcommand\listtablename{List of Tables}
\fi
\ifdefined\figurename
  \renewcommand*\figurename{Figure}
\else
  \newcommand\figurename{Figure}
\fi
\ifdefined\tablename
  \renewcommand*\tablename{Table}
\else
  \newcommand\tablename{Table}
\fi
}
\@ifpackageloaded{float}{}{\usepackage{float}}
\floatstyle{ruled}
\@ifundefined{c@chapter}{\newfloat{codelisting}{h}{lop}}{\newfloat{codelisting}{h}{lop}[chapter]}
\floatname{codelisting}{Listing}
\newcommand*\listoflistings{\listof{codelisting}{List of Listings}}
\makeatother
\makeatletter
\makeatother
\makeatletter
\@ifpackageloaded{caption}{}{\usepackage{caption}}
\@ifpackageloaded{subcaption}{}{\usepackage{subcaption}}
\makeatother
\usepackage{bookmark}
\IfFileExists{xurl.sty}{\usepackage{xurl}}{} % add URL line breaks if available
\urlstyle{same}
\hypersetup{
  pdftitle={aMindLikeARazor},
  pdfauthor={Gabriel X. Martinez},
  colorlinks=true,
  linkcolor={blue},
  filecolor={Maroon},
  citecolor={Blue},
  urlcolor={Blue},
  pdfcreator={LaTeX via pandoc}}


\title{aMindLikeARazor}
\author{Gabriel X. Martinez}
\date{2026-11-21}
\begin{document}
\maketitle

\renewcommand*\contentsname{Table of contents}
{
\setcounter{tocdepth}{2}
\tableofcontents
}

\bookmarksetup{startatroot}

\chapter*{Preface}\label{preface}
\addcontentsline{toc}{chapter}{Preface}

\markboth{Preface}{Preface}

The CEO of a popular razor-making company is being interviewed by a
famous business magazine. ``Is it hard to recruit the best minds to work
for you?'' asks the interviewer. ``Well,'' says the CEO, ``our people
could have worked for NASA or Ferrari. But when we tell them that they
will be working on the interaction between living biological tissue and
the sharpest object known to man, we have no problem getting them to
sign up.''

Economics is about the interaction between the sharpest known object --
your razor-like intelligence, the sharpest logic, the glint of
mathematical steel -- and the living, breathing, pulsating ordinary
business of life. We answer questions like: what makes some countries
rich and others poor? Why do we fall into high unemployment, and how do
we get out of it? Why do financial crises always seem to catch us by
surprise? Why are things so expensive and who's to blame and who's in
charge? How does our country interact with the rest of the world? What
determines prices and why do they change? Why does the free market work
so well so much of the time, delivering trillions of dollars of goods
and services without a glitch and without a massive coordinating plan?
How do monopoly power or information failure mess up that beautiful
thing? How do we make it work when what you do affects me? What does
economics have to say about seemingly uneconomic subjects, like the
family or people's work?

This is an Economics book. This means that you should be ready to work
hard, and have fun working hard. It's written so that introductory
students can get all they need -- and if they want to learn more, the
intermediate and advanced topics are right there. And it is written so
that intermediate-level students can easily go back to the topics they
(of course) learned and almost remember but just need a refresher for.
It's written for beginning students who remember their Algebra II, for
intermediate students who loved their Pre-Calculus, and for more
advanced students who think better when they take derivatives and
maximize.

To learn more about Quarto books visit
\url{https://quarto.org/docs/books}.

\begin{Shaded}
\begin{Highlighting}[]
\DecValTok{1} \SpecialCharTok{+} \DecValTok{1}
\end{Highlighting}
\end{Shaded}

\begin{verbatim}
[1] 2
\end{verbatim}

\bookmarksetup{startatroot}

\chapter*{Introduction}\label{introduction}
\addcontentsline{toc}{chapter}{Introduction}

\markboth{Introduction}{Introduction}

This is a book created from markdown and executable code.

See Knuth (1984) for additional discussion of literate programming.

\begin{Shaded}
\begin{Highlighting}[]
\DecValTok{1} \SpecialCharTok{+} \DecValTok{1}
\end{Highlighting}
\end{Shaded}

\begin{verbatim}
[1] 2
\end{verbatim}

\bookmarksetup{startatroot}

\chapter{The Macroeconomic
Perspective}\label{the-macroeconomic-perspective}

You work hard, develop useful skills, and make careful financial
decisions.~ What important things about your economic future are still
outside your individual control? The opportunities to earn may be
abundant or scarce. Your earnings may be able to buy all you need or
leave you short. The economy might just not have that much to offer, or
it might be plentiful.

What are the conditions that shape those opportunities? How do our
decisions collectively help create those conditions?

\section{Reading}\label{reading}

IMF / \emph{Finance \& Development},
``\href{https://www.imf.org/external/pubs/ft/fandd/basics/14_gdp.htm}{Gross
Domestic Product: An Economy's All}.'' What is one reason to care about
GDP growth, and one reason GDP cannot tell us everything we want to
know?

Vermann, Katarina (2012)
\href{https://www.stlouisfed.org/publications/page-one-economics/2012/05/01/wait-is-saving-good-or-bad-the-paradox-of-thrift}{Wait,
Is Saving Good or Bad? The Paradox of Thrift}

\textbf{Microeconomics} is about individual decisions.
\textbf{Macroeconomics} is about how those individual decisions lead to
outcomes for all of us. \textbf{Micro} individual decisions are made
with constraints and in response to incentives. \textbf{Macro} context
is about how those constraints come about and how those incentives
interact.

The Macro perspective is about feedback and interaction, which means
that it is about the big picture. Your spending is my income. If we all
make the same (perfectly rational) individual decision, the result might
be bad for all of us.

In Macroeconomics, we focus on five core questions: What determines how
well people can live and why do living standards rise? Why do some
economies grow faster than others? Why is it so hard (or easy) for
millions of people to find work? What drives changes in the purchasing
power of my money? Can society, through its government, come up with
policies to manage the economy?

These are really important questions. What opportunities are available
to me? What opportunities will be available to my children? How much of
my job outlook depends on where and when I am looking? Will next year's
raise actually let me buy more? Why does the economic situation change,
and can we do something about it?

\section{Three Graphs}\label{three-graphs}

The blue, solid line is the actual size of the economy of the United
States since 1947. Notice the periods when the economy shrank, but also
notice the long-term pattern of growth around Potential Output, the
dashed red line.

\begin{figure}

\centering{

}

\caption{\label{fig-realgdp-potgpd}}

\end{figure}%

How hard is it to find work? If you are looking for work, and can't find
it, you are considered ``unemployed.'' There are many reasons for
unemployment, some of which depend on the state of the economy (which
explains the difference between the blue solid line and the red dashed
line).

\begin{figure}

\centering{

}

\caption{\label{fig-unrate-nrou}}

\end{figure}%

How quickly are things getting more expensive? That depends on which
``things'' we are talking about. The blue solid line shows you the
\emph{growth rate} of the prices of everything consumers buy. The red
dashed line is also a growth rate of prices, but excluding shocks that
are quickly reversed.

\begin{figure}

\centering{

}

\caption{\label{fig-pcepi-pcepilfe}}

\end{figure}%

\section{The Habits of the
Macroeconomist}\label{the-habits-of-the-macroeconomist}

To think like a macroeconomist is to

\begin{itemize}
\item
  \textbf{Follow Connections}. If you spend less at the local
  restaurant, what happens to the restaurant workers' income? Everyone's
  spending is someone else's income; people save so that others may
  invest in new capital; pumping money into the economy increases
  spending but also inflation.
\item
  \textbf{Add It Up}. If \emph{everyone} spends less, what happens to
  everyone's income? In Micro, one person's spending doesn't change the
  entire economy. In Macro, we \emph{aggregate}, meaning that we add up:
  add up spending on all products, add up the incomes of workers and
  owners, add up the production of farmers, factories, and financial
  firms, add up everything that can be considered money.
\item
  \textbf{Keep It Real}. It's not about the dollars you get paid -- it's
  about how much they can buy. It's about how living standards change,
  how your material welfare changes. If your wages rise by 10\% but
  prices rise by 11\%, you are no richer. If a country produces 20 times
  more than another country but has 30 times more people, it's not more
  prosperous.
\item
  \textbf{Separate the Long and Short of It}. A country is richer
  because of a long-term trend, as its economy grew over generations.
  But its people might have trouble finding work in any particular year
  due to temporary fluctuations around that trend. The distinction
  between short-run cycles and the long-run path is crucial to
  macroeconomics.
\item
  \textbf{Cut it some slack.} On any given day, you may be feeling a
  little low or very energized, operating below your potential or above
  it -- and you always need to have some rest. An economy can operate
  above its potential, or shot of it. If the economy is producing less
  than its capacity, you see unemployment. Inflation typically picks up
  if the economy is operating above capacity.
\item
  \textbf{Look at the forest.} ``How is the economy doing?'' is not the
  kind of question you can ask by just looking at one number. You need
  to look at the whole picture: GDP growth, inflation, and unemployment
  are the three big statistics, Even an accurate aggregate can conceal
  differences across households and communities.
\item
  \textbf{Explain, Test, Repeat.} It's easy to tell a simple, memorable
  story. But the truth is often more interesting and more complex. Build
  a real explanation of the world around you, and then test it by
  asking, ``if my explanation is true, then what would I observe?''
\end{itemize}

\section{Interpret the Data}\label{interpret-the-data}

Look back at Figure~\ref{fig-realgdp-potgpd}. Can output rise over a
long time but fall over a shorter period?

Search the internet for this answer: if unemployment is 5\%, does that
mean that 95\% of the population is employed? Look at
Figure~\ref{fig-unrate-nrou} for last month's unemployment rate.

Figure~\ref{fig-pcepi-pcepilfe} shows that between July 2022 and
December 2024, inflation fell from more than 7\% to 2.5\%. Does this
mean that prices returned to where they started? Hint: the definition of
inflation is ``the growth rate of prices.''

Macroeconomists love graphs. But not because they really explain
anything. Graphs tell us what happened, but it is up to a good theory to
explain \emph{why} and \emph{how} it happened. We need a causal
mechanism and complementary evidence.

We also love variables like ``GDP per capita,'' which is the size of the
economy divided by the population. But we know that a single variable
can never be a complete answer to complicated questions like ``are
people better off?'' So we want to \textbf{measure} economic
performance, and we want to \textbf{explain} it, but we also want to
\textbf{evaluate} it. Never confuse ``up'' with ``good.''

\section{Think About It}\label{think-about-it}

Give one example of how an individually sensible decision could have a
different effect when many people make it.

Choose one dashboard measure, like GDP, unemployment, or inflation. What
does it tell us, and what does it not tell us?''

\bookmarksetup{startatroot}

\chapter{Summary}\label{summary}

In summary, this book has no content whatsoever.

\begin{Shaded}
\begin{Highlighting}[]
\DecValTok{1} \SpecialCharTok{+} \DecValTok{1}
\end{Highlighting}
\end{Shaded}

\begin{verbatim}
[1] 2
\end{verbatim}

\bookmarksetup{startatroot}

\chapter*{References}\label{references}
\addcontentsline{toc}{chapter}{References}

\markboth{References}{References}

\protect\phantomsection\label{refs}
\begin{CSLReferences}{1}{1}
\bibitem[\citeproctext]{ref-knuth84}
Knuth, Donald E. 1984. {``Literate Programming.''} \emph{Comput. J.}
(USA) 27 (2): 97--111. \url{https://doi.org/10.1093/comjnl/27.2.97}.

\end{CSLReferences}




\end{document}
